% Ad Familiares 2.16 --- headnote
% ad-familiares-02-16, 4 May 49 BC, Cumae

\textit{Cicero to M. Caelius Rufus, written from his villa
at Cumae on the fourth day before the Nones of May ---
Perseus dateline \textit{Scr. in Cumano iv Non. Mai. a. 705
(49)}, that is, 4 May 49 BC. The salutation styles Cicero
\textit{imperator}, the title that has clung to him since his
Cilician proconsulship and the lictors of which (the ``awkward
pomp of my lictors'') he complains in the body of the letter.
Caelius, who had been Cicero's quaestor of sorts at Rome and
his lively correspondent throughout the Cilician year, has
gone over to Caesar; he has written a coaxing, half-anxious
letter urging Cicero not to take ship for Greece and warning
him of risks to himself and his family if he does. Cicero
replies a week after his letter to Servius Sulpicius (Fam.
4.2) and is on the brink of the very decision he denies to
Caelius he has made.}

\textit{The letter is among the most exposed psychological
documents of the spring of 49: Cicero protests too much that
he has no intention of joining Pompey, that he is merely
looking for a quiet \textit{solitudo} in Italy, that his
fondness for the seaside villas means nothing --- while
admitting, in the same breath, that he would gladly take ship
``for the sake of quiet,'' that the dignity of his imperatorial
laurel has become a target for malice, and that his model in
neutrality is the late Q. Hortensius. The reference to
Dolabella, Cicero's son-in-law and at this point a young
Caesarian, is the closest the letter comes to confession:
amid the wreckage Cicero had hoped that Dolabella's debts
might at least have been wiped out by his change of side, and
his ``dishonourable days'' for his father-in-law are the
weeks when he served Caesar in the city while Cicero was
still wavering. The closing throwaway about the gifted
\textit{toga praetexta} for the boy of Oppius --- one of
Caesar's chief agents at Rome --- is the carefully placed
proof that he has not yet broken with the other side.}
